\subsection{Tournament Performance Model}
\label{sec:strategy_tournament}

The first two identification strategies measure the effect of lucky loser entry on \emph{ranking}---an aggregate outcome that conflates match wins, opponent quality, tournament tier, and the mechanical accumulation of ranking points. A player who goes 2W--1L at a Grand Slam earns far more ranking points than one who goes 2W--1L at an ATP 250, yet both register identically in a count-based outcome variable. To ask whether LL entry makes a player more likely to win matches against specific opponents in specific competitive contexts, I need an outcome that preserves the full structure of tournament competition.

\subsubsection{Defining Tournament Performance}

Consider player $i$ who is in the LL candidate pool at qualifying event $e$ (the reference event, $t = 0$). At each subsequent tournament $t = 1, 2, \ldots$, the player faces a sequence of opponents $\{j_1, j_2, \ldots\}$ determined by the draw. I define the player's \textit{tournament performance} $F_{it}$ as their complete trajectory at that event: the ordered sequence of match outcomes against identified opponents. For example, if the player beats $j_1$ and $j_2$ but loses to $j_3$, then $F_{it} = (W, W, L)$ against the opponent set $\{j_1, j_2, j_3\}$.

The probability of observing a specific trajectory is the product of the constituent match probabilities. In the example above:
\begin{equation}
\label{eq:tournament_performance}
P\bigl(F_{it} = \text{2W--1L} \mid \{j_1, j_2, j_3\}\bigr) = P(W_{it1} = 1 \mid j_1) \times P(W_{it2} = 1 \mid j_2) \times P(W_{it3} = 0 \mid j_3)
\end{equation}
More generally, for a trajectory of $w$ wins followed by a loss:
\begin{equation}
\label{eq:sequential_likelihood}
P(F_{it}) = \left(\prod_{m=1}^{w} p_{itm}\right) \cdot \left(1 - p_{it,w+1}\right)
\end{equation}
where $p_{itm}$ is the probability that player $i$ wins match $m$ at tournament $t$, and the final loss term is absent if the player wins the tournament outright. This is an inhomogeneous geometric distribution: a sequence of independent Bernoulli trials with match-specific success probabilities, stopped at the first failure.

The advantage of this formulation over a scalar summary---win count, ranking points earned, or an indicator for reaching a given round---is that it preserves opponent identity and match-level covariates at every round. The same record of 2W--1L has a very different likelihood depending on whether the three opponents were seeds or qualifiers, whether the surface favors the focal player, and how the player's pre-treatment strength compares to each opponent.

\subsubsection{The Match-Level Model}

Each match probability $p_{itm}$ is modeled as a logistic function of match-level covariates, pre-treatment characteristics, the LL treatment indicator from event $e$, and the control function correction:
\begin{equation}
\label{eq:match_logit}
p_{itm} = \Lambda\bigl(\beta' X_{ijm} + \gamma' Z^{pre}_{ie} + \delta \cdot D_{ie} + \rho \cdot \hat{v}_{ie}\bigr)
\end{equation}
where $\Lambda(\cdot)$ is the logistic CDF, $D_{ie}$ indicates whether player $i$ received LL entry at reference event $e$, and $\hat{v}_{ie}$ is the generalized residual computed analytically from the selection probability $P_i^{LL}$ (Section~\ref{sec:strategy_iv}). The match-level pairwise covariates $X_{ijm}$ include ranking points difference / 1000 (linear and quadratic), Elo difference / 100 (linear and quadratic), surface Elo difference / 100 (linear and quadratic), a Laplace-smoothed head-to-head win proportion, the number of prior head-to-head encounters, the age difference, and surface indicators (clay and grass)---the same logit specification used for the selection probability $P_i^{LL}$ and the treatment dose analysis. The pre-treatment covariates $Z^{pre}_{ie}$ include ranking points / 1000 (linear and quadratic), Elo / 100 (linear and quadratic), surface Elo / 100 (linear and quadratic), and player age, all measured at the date of the qualifying loss at event $e$. The quadratic terms in both the pairwise and pre-treatment covariates allow for nonlinear returns to ranking, ability, and surface specialization.

Because the log-likelihood of the sequential product (\ref{eq:sequential_likelihood}) decomposes additively across matches---$\sum_m \bigl[w_{itm} \log p_{itm} + (1 - w_{itm}) \log(1 - p_{itm})\bigr]$---estimation reduces to a single logit on the stacked match-level dataset. No distributional assumptions beyond the logit link are required; the sequential tournament structure emerges from the product of match-level contributions.

The parameter $\delta$ answers the question: conditional on pre-treatment ability and the specific opponent faced, does having received an LL slot at event $e$ increase the probability of winning matches at future tournaments $t > 0$? The control function term $\hat{v}_{ie}$ absorbs the correlation between LL receipt and unobserved player quality. Its coefficient $\rho$ directly tests for endogeneity: if $\rho \neq 0$, the na\"ive logit without the correction is inconsistent.

\subsubsection{Identification and the Control Function}

At non-Grand-Slam events, identification relies on the control function approach described in Section~\ref{sec:strategy_iv}. The selection probability $P_i^{LL}$ is constructed from logit-predicted match-winning probabilities combined via exact enumeration of qualifying match outcomes (Section~\ref{sec:strategy_iv}). The generalized residual $\hat{v}_{ie}$ is computed analytically from $P_i^{LL}$ and enters the match-level logit alongside the treatment indicator $D_{ie}$. Tournament level and playing surface do not enter the computation of $P_i^{LL}$ because LL allocation at non-Grand-Slam events depends on the ranking ordering among qualifying losers and the stochastic withdrawals of main draw players---not on the event's surface or tier. These variables enter the outcome equation as match-environment controls affecting $p_{itm}$ but not the selection probability.

\subsubsection{Pre-Treatment Covariates}

If LL entry at event $e$ causally affects the player's Elo rating---through main draw match results, ranking points, and the tournament access they generate---then conditioning on post-treatment Elo at tournament $t > 0$ absorbs part of the causal pathway from $D_{ie}$ to $F_{it}$. I address this by freezing the focal player's covariates at the date of the qualifying loss at event $e$: the Elo rating, surface Elo, and head-to-head record in $X_{itm}$ reflect the player's ability at the moment just before treatment could take effect. The opponent's Elo is measured at the match date, since it is unaffected by the focal player's treatment status. This asymmetric differencing---pre-treatment focal ability minus contemporaneous opponent strength---identifies how the LL experience shifts win probability beyond what pre-existing ability would predict.

\subsubsection{Event-Cohort Structure}

Each qualifying loss defines an event $e$. The player is tracked over subsequent main draw tournaments within a window defined as the minimum of 10 tournaments, 365 days, and the date of the player's next qualifying loss. The last constraint---truncation at the next LL candidate event---ensures that within each cohort, the only relevant treatment is $D_{ie}$ from the originating event, preventing contamination from overlapping treatment spells. A player who appears in the LL candidate pool at multiple events generates multiple cohorts, potentially with different treatment statuses. Standard errors are obtained via player-level block bootstrap that resamples entire player histories, preserving within-player correlation across events and matches.

\subsubsection{Compounding: From Match-Level to Tournament-Level Effects}

The parameter $\delta$ is denominated in log-odds of winning a single match. Because tournament performance $F_{it}$ is a \textit{product} of match probabilities, the per-match effect compounds across rounds. To translate $\delta$ into tournament-level impacts, I compute three derived quantities using counterfactual simulation over the observed data.

The \textit{match-level average marginal effect} is the sample mean of $p_{itm}(D=1) - p_{itm}(D=0)$ across all observed matches, where $p_{itm}(D=x) = \Lambda(X_{itm}'\hat{\beta} + \hat{\delta} \cdot x)$. This gives the average per-match probability shift, integrated over the actual distribution of opponents and competitive contexts.

The \textit{tournament-level shift in expected wins} is computed for each player-tournament as $\Delta E[W_{it}] = \sum_{r=1}^{R}\prod_{s=1}^{r} p_{its}(1) - \sum_{r=1}^{R}\prod_{s=1}^{r} p_{its}(0)$, using the identity $E[W] = \sum_r P(W \geq r)$. This captures how many additional rounds the player is expected to survive at tournament $t$ due to having received the LL slot at event $e$.

The \textit{shift in tournament win probability} is $\pi_{it}(1) - \pi_{it}(0)$, where $\pi_{it}(x) = \prod_{r=1}^{R} p_{itr}(x)$. Because tournament win probabilities are typically small, the risk ratio $\pi_{it}(1)/\pi_{it}(0)$ is often more informative. A modest per-match $\delta$ can produce a substantial tournament-level effect through multiplicative compounding: for instance, a 3 percentage point match-level advantage compounds to a roughly 50\% increase in tournament win probability over five rounds.

All three quantities are computed for every player-tournament observation and then averaged, integrating over the actual draw compositions and opponent strengths in the data rather than relying on hypothetical scenarios.
